\section{Two rapidly degenerating ghost paths}
\label{app:raw-tail-stress}

These two coefficient paths illustrate why the complete physical gain
must be estimated before its individual factors are simplified. In the
quadratic chart, large common logarithms cancel in the exact centre
matching. In the compact cubic chart, the positive bottleneck forces the
negative lip extremely close to its simple root, and the resulting
endpoint power remains in the output gain.
Both paths use the affine input and divided closing of
Lemma~\ref{lem:closing}, on the original buffered carrier below a fixed
positive raw cap. Their actual centre maps and completed gain quotient
are those of Proposition~\ref{rawpair:prop-original-tails}.
We denote the circular constant by
\(\pi_{\rm circ}\) when a coefficient \(\pi_0\) occurs in the same formula.

\subsection{A quadratic path with an exponentially small chart ratio}
\label{rawstress:subsec-quadratic}

Fix \(\pi_0\) strictly inside the retained ghost band at \(k=0\), so
\[
 \pi_\flat<\pi_0<\sqrt{12},\quad \pi_0\le\Pi,\qquad
                  -\Delta_0<\Delta(0,\pi_0)<0.
\]
Retain a compact raw interval
\(\mathcal A\Subset[d_0,a_0)\), where \(a_0\) is the original positive
outward reserve of Corollary~\ref{deg:cor-positive-large-S-outer}.

\begin{proposition}\label{rawstress:prop-quadratic}
For \(g\downarrow0\), take
\begin{equation}\label{rawstress:eq-quadratic-coefficients}
 \begin{gathered}
 k=e^{-1/g^2},\qquad A=1,\quad C=1+g,\quad \mu_3=-g-g^2,\\
 \mu_2=(3+\pi_0k)g^2,\qquad
 \mu_1=-(2+\pi_0k+k^2)g^3.
 \end{gathered}
\end{equation}
For every sufficiently small \(g>0\), the lips with
\(a\in\mathcal A\), \(u=a/(gk)\), give actual inputs of the original
carrier. Their matched negative radii \(R=kr\) satisfy, uniformly on
\(\mathcal A\),
\begin{equation}\label{rawstress:eq-quadratic-radius}
 R\longrightarrow R_{\pi_0}:=
 \frac{3e^{3D_{\pi_0}}}{1+e^{3D_{\pi_0}}},\qquad
 D_{\pi_0}=\frac{\pi_0\pi_{\rm circ}}
                       {3\sqrt{12-\pi_0^2}}.
\end{equation}
Their raw central heights obey
\begin{equation}\label{rawstress:eq-quadratic-height}
 \log H(a)=-\frac1{3g^3}-\frac2{3g^2}
                                    +O(g^{-1}+|\log g|).
\end{equation}
The actual finite channel has \(K_{\rm fin}>0\) and \(m\to0\).
The completed physical derivative satisfies
\(0<\Pi'\le C(d_0)g^3\) on these inputs.
\end{proposition}

\begin{proof}
The coefficients tend to \(\theta_0\), and the small transverse zero is
exactly \(\gamma=g\). Substitution in \eqref{eq:polynomials} gives
\[
 P_\theta(g+t)=-M(t),\qquad
 M(t)=t^3+3gt^2-\pi_0kg^2t+k^2g^3.
\]
Thus these are fibres of \eqref{deg:quadratic-field}. The finite
discriminant expansion in Section~\ref{sec:normalization} gives
\[
 K_{\rm fin}=C-A+O(|\theta-\theta_0|^2)=g+O(g^2)>0.
\]
The actual multiplier \eqref{eq:finite-multiplier} consequently tends
to zero and eventually satisfies the last multiplier restriction.

Put \(S=ku=a/g\). Then \(u\ge2M\), \(S\ge\delta\), and \(a<a_0\)
for sufficiently small \(g\). Corollary~\ref{deg:cor-positive-large-S-outer}
places the positive raw image in the prescribed buffer. Its reserve is
fixed using the raw section margin, with \(a_0<d_q/L\); it applies when
\(a\) is fixed while \(S\to\infty\). The inward construction in
\eqref{bp-tail:positive-inward-variation}--\eqref{bp-tail:positive-inward-absolute}
gives a positive actual centre height. Its monotonicity and the original
comparison give the unique negative match
\eqref{bp-tail:whole-tail-matching}. Corollary
\ref{deg:cor-negative-absolute-outer} constructs its outward continuation
in the negative buffer, including the original bridge. The two inward
graphs have the same height at \(v=g\); scalar uniqueness joins them
to one through orbit. The finite closing is therefore evaluated on its
prescribed image. The inward estimate on the fixed coefficient band gives
\(H\le(gk)^2Cg e^{-c/g}\), and the raw speed at the centre is
\(k^2g^3+\ell(g)H\). Both tend to zero. The orbits therefore satisfy the fixed height cap and speed restriction;
their positive height places them above the stopped endpoint.

To determine the matched radius, we keep the coefficient-dependent
logarithms in the two action primitives until their difference is taken. Set
\[
 c=\mathfrak c,\quad q=c^{-1},\quad
 \omega=\sqrt{-\Delta},\qquad
 \mathcal Q_+(\xi)=(\xi-d)^2+\omega^2,\quad
 \mathcal Q_-(r)=(r+d)^2+\omega^2.
\]
Then \(P(\xi)=(c+k\xi)\mathcal Q_+(\xi)\) and
\(P_-(r)=(c-kr)\mathcal Q_-(r)\). Define
\begin{equation}\label{rawstress:eq-quadratic-partial-coefficients}
 \beta_q=\frac{1+2g+g(2dk+k^2/c^2)}{c+2dk+k^2/c^2},\qquad
 \rho_q=k(g-\beta_q),\qquad \chi_q=-\rho_q/c^2.
\end{equation}
Direct multiplication yields
\begin{equation}\label{rawstress:eq-quadratic-partial-fractions}
 \begin{aligned}
 \frac{\xi(1+2g+gk\xi)}{P(\xi)}
   &=\frac{\rho_q}{c+k\xi}
                      +\frac{\beta_q\xi+\chi_q}{\mathcal Q_+(\xi)},\\
 \frac{r(1+2g-gkr)}{P_-(r)}
   &=-\frac{\rho_q}{c-kr}
                      +\frac{\beta_qr-\chi_q}{\mathcal Q_-(r)}.
 \end{aligned}
\end{equation}
In particular,
\[
 \beta_q=\frac{1+2g}{3}+O(k),\qquad
 \rho_q/k=g-\beta_q=-\frac{1-g}{3}+O(k),\qquad
                         \rho_q/k+\beta_q=g.
\]
The last identity is exact.

Let \(\mathcal A_+(u)=\int_0^u\xi B_+/P\,d\xi\) and
\(\mathcal A_-(r)=\int_0^r vB_-/P_-\,dv\), with the same original
centre as base point. Integrating \eqref{rawstress:eq-quadratic-partial-fractions}
gives
\begin{align}
 \mathcal A_+(u)={}&(g-\beta_q)\log\frac{c+ku}{c}
       +\frac{\beta_q}{2}\log\frac{\mathcal Q_+(u)}q\notag\\
 &+\frac{\chi_q+\beta_qd}{\omega}
   \left(\arctan\frac{u-d}{\omega}+\arctan\frac d\omega\right),\notag\\
 \mathcal A_-(r)={}&(g-\beta_q)\log\frac{c-kr}{c}
       +\frac{\beta_q}{2}\log\frac{\mathcal Q_-(r)}q\notag\\
 &-\frac{\chi_q+\beta_qd}{\omega}
   \left(\arctan\frac{r+d}{\omega}-\arctan\frac d\omega\right).
 \label{rawstress:eq-quadratic-primitives}
\end{align}
For \(u=S/k\), \(r=R/k\), the common logarithmic term is exactly
\(\beta_q\log(1/k)\). Subtracting the two primitives gives
\begin{align}
 \mathcal A_+(S/k)-\mathcal A_-(R/k)={}&
  (g-\beta_q)\log\frac{c+S}{c-R}\notag\\
 &+\frac{\beta_q}{2}\log
   \frac{(S-kd)^2+k^2\omega^2}{(R+kd)^2+k^2\omega^2}\notag\\
 &+\frac{\chi_q+\beta_qd}{\omega}
  \left(\arctan\frac{S-kd}{k\omega}
             +\arctan\frac{R+kd}{k\omega}\right).
 \label{rawstress:eq-quadratic-action-difference}
\end{align}
The coefficient of the remaining large \(\log S\) in the combined
radial terms is \(g\).

The two inward constructions give the actual centre heights
\[
 \eta_0^+=g\ell_0P(u)e^{-\mathcal A_+(u)/g}E_+,\qquad
 \eta_0^-=g\ell_0P_-(r)e^{-\mathcal A_-(r)/g}E_-,
 \qquad E_\pm,E_\pm^{-1}\le C.
\]
The bounded factors contain the nonlinear ratio integrals and the
compact continuations to the centre. Their bounds follow from the action-clock decay
in the two inward constructions. Matching these actual heights gives
\begin{equation}\label{rawstress:eq-quadratic-exact-memory}
 \mathcal A_+(u)-\mathcal A_-(r)
    =g\log\frac{P(u)}{P_-(r)}+g\log\frac{E_+}{E_-}.
\end{equation}
The initial polynomial factors in these heights have the exact quotient:
\begin{equation}\label{rawstress:eq-quadratic-polynomial-memory}
 \frac{P(S/k)}{P_-(R/k)}
   =\frac{k^2-\pi_0kS+3S^2+S^3}
          {k^2+\pi_0kR+3R^2-R^3}.
\end{equation}
On a fixed compact subinterval of \(0<R<3\), its logarithm is
\(3\log S-\log[R^2(3-R)]+O(S^{-1}+k)\).
Thus its contribution to \eqref{rawstress:eq-quadratic-exact-memory}
is \(O(g(1+|\log g|))\). The common terms containing \(\log k\) cancel in the exact action and
polynomial identities before the estimate is taken, although
\(g\log(1/k)=g^{-1}\to\infty\) on this path.

For a fixed \(R\in(0,3)\), \eqref{rawstress:eq-quadratic-action-difference}
tends to
\[
 D_{\pi_0}+\frac13\log\frac{3-R}{R}.
\]
Indeed \(d\to\pi_0/6\),
\(\omega\to\sqrt{12-\pi_0^2}/6\), and both arctangents tend to
\(\pi_{\rm circ}/2\). This limiting expression has the single zero
\(R_{\pi_0}\) in \eqref{rawstress:eq-quadratic-radius}.
Choose fixed \(R_{\rm lo}<R_{\pi_0}<R_{\rm hi}\) in \((0,3)\).
The exact matching residual has opposite signs at these two radii for
small \(g\), uniformly on \(\mathcal A\). Strict monotonicity of the
actual negative centre map places its match between them. The
expansion is uniform on that compact interval; shrinking the brackets
proves \eqref{rawstress:eq-quadratic-radius}. The limiting gap from the
negative root is positive for this fixed \(\pi_0\).

The positive primitive also gives
\[
 \mathcal A_+(u)=\frac{1+2g}{3}\log(1/k)+g\log S+O(1).
\]
Since \(H=\ell_0M(a)e^{-\mathcal A_+(u)/g}E_+\) and
\(M(a)\to a^3\), this proves \eqref{rawstress:eq-quadratic-height}.
In the scaled centre coordinate,
\[
 \log\eta_0^+=-\frac1{3g^3}+\frac4{3g^2}
                                      +O(g^{-1}+|\log g|).
\]

Finally retain the two complete exterior gains. With
\(\nu=6g/(1+2g)\), \(\lambda_g=2\nu\),
\(p=\widehat P(R)/R^2\), and \(b_R=-R^2/2\), the negative gain is
\begin{equation}\label{rawstress:eq-quadratic-full-gain}
 \Gamma_{-,b_R}'=
 c_-g^2\ell_0^{-1-\nu}g^{-\nu}(gR)^{-\lambda_g}
                        p^{-\nu}\mathcal K_-,\qquad
                         0<\mathcal K_-\le C.
\end{equation}
Its original label change is
\(\Gamma_{-,b_r}'=k^2\Gamma_{-,b_R}'\), where \(b_r=-r^2/2\).
The positive input gain in \(b_u=-u^2/2\) retains \((gk)^2\).
Using the actual centre variations gives
\begin{equation}\label{rawstress:eq-quadratic-consumed-gains}
 \begin{gathered}
 0<\frac{d\Gamma_-}{d\log\eta_0^-}
 \le Cc_-g^3g^{-\nu-\lambda_g}
                       R^{2-\lambda_g}p^{1-\nu}\le Cc_-g^3,\\
 \frac{d\Gamma_+}{d\log\eta_0^+}\ge c(d_0)c_+a^3.
 \end{gathered}
\end{equation}
The bounded negative factor follows from the upper density estimate
\eqref{bp-tail:negative-density-value}. Its root dwell has nonpositive
density, as shown in \eqref{bp-tail:negative-exact-density}; it does
not require an inverse bound on \(\mathcal K_-\). The exact quotient
of \eqref{rawstress:eq-quadratic-consumed-gains} is the completed
physical \(\Pi'\). This proves its stated bound. If the exterior gain is factored differently, the term
\(k^{-\lambda_g}=\exp(12/[g(1+2g)])\) is unbounded on this path;
the exact label change and centre variation account for it before the
quotient is estimated.
\end{proof}

\subsection{A compact cubic path with a very small positive bottleneck}
\label{rawstress:subsec-cubic}

We write \(\tau_{\rm c}=\tau_{\rm cub}\) for the raw cubic scale and
\(\kappa_{\rm c}\) for its positive ghost gap. Choose a compact block
with \(U_0\ge1\) and \(G_0>1/3\). For small \(\kappa_{\rm c}>0\), set
\begin{equation}\label{rawstress:eq-cubic-coefficients}
 \begin{gathered}
 \kappa_{\rm c}=e^{-1/\tau_{\rm c}^2},\qquad
 c_{\rm c}=(1+\kappa_{\rm c}^2)^{-1},\qquad \alpha_{\rm c}=1,\\
 U_{\rm c}=2c_{\rm c}-c_{\rm c}^{-1},\qquad
 G_{\rm c}=(c_{\rm c}-2)/3.
 \end{gathered}
\end{equation}
These coefficients belong to the block for sufficiently small
\(\tau_{\rm c}\), and \(U_{\rm c}\ge0\). Their axis polynomial is
\begin{equation}\label{rawstress:eq-cubic-polynomial}
 P_{\rm c}(v)=(v+c_{\rm c})
                   \{(v-1)^2+\kappa_{\rm c}^2\},
 \qquad \Delta_{\rm c}=-\kappa_{\rm c}^2<0.
\end{equation}

There is a fixed raw reserve \(a_{\rm res}>0\) inside the positive
window for this path. To see this, at \(\tau_{\rm c}=0\) consider the
positive raw graph with \(M(t)=t^3\), initial height
\(z(a)=\ell_0a^3\), and the organizing passive coefficients.
The near-lip comparison of Lemma~\ref{deg:lem-raw-growing-section}
puts \(z/t^2\) in a fixed positive strip at \(t=2a\). The subsequent
equation is
\(t(z/t^2)'=1-2z/t^2+O(t)\), with a negative normal diagonal on
that strip. Its fixed-port value converges to the organizing strong
graph as \(a\downarrow0\). The strict window margin therefore
admits all these limiting raw images for \(0<a\le a_{\rm res}\),
after decreasing \(a_{\rm res}\). Ordinary dependence on the regular
raw interval from a fixed \(a>0\) to the port preserves this margin
for small \(\tau_{\rm c}\), uniformly on compact raw subintervals.
Choose the auxiliary \(d_0\), after the original windows and reserves,
so that \(2d_0<a_{\rm res}\), while retaining the small-error
requirements of its construction. We may therefore use \(a\in[d_0,2d_0]\) within the originally
prescribed physical carrier.

\begin{proposition}\label{rawstress:prop-cubic}
The path \eqref{rawstress:eq-cubic-coefficients} is realized in the full
coefficient germ by
\begin{equation}\label{rawstress:eq-cubic-full-coefficients}
 \begin{gathered}
 A=1,\quad C=1+\tau_{\rm c},\quad
 \gamma=\tau_{\rm c}G_{\rm c},\quad
 \mu_3=-\gamma-\gamma^2,\\
 \mu_2=\tau_{\rm c}^2(U_{\rm c}+3G_{\rm c}^2),\qquad
 \mu_1=-\tau_{\rm c}^3(1+U_{\rm c}G_{\rm c}+2G_{\rm c}^3).
 \end{gathered}
\end{equation}
For sufficiently small \(\tau_{\rm c}>0\), the lips with
\(a\in[d_0,2d_0]\), \(u=a/\tau_{\rm c}\), are original admitted
inputs. If \(X=c_{\rm c}-r\) is their matched negative root distance,
then uniformly on this interval
\begin{equation}\label{rawstress:eq-cubic-root-distance}
 \kappa_{\rm c}\log(1/X)\longrightarrow2\pi_{\rm circ}.
\end{equation}
Their raw central height satisfies
\begin{equation}\label{rawstress:eq-cubic-height}
 \log H(a)=
 -\frac{\pi_{\rm circ}B(1)}
             {(1+c_{\rm c})\tau_{\rm c}\kappa_{\rm c}}
                  +O(\tau_{\rm c}^{-1}+|\log\tau_{\rm c}|),
\end{equation}
so \(\tau_{\rm c}\kappa_{\rm c}\log H(a)\to-\pi_{\rm circ}/2\).
The finite channel has \(K_{\rm fin}>0\), \(m\to0\), and the full
physical derivative obeys \(0<\Pi'\le C(d_0)\tau_{\rm c}^3\).
\end{proposition}

\begin{proof}
Translation of the raw polynomial gives
\[
 P_\theta(\gamma+t)=-M_{\rm c}(t),\qquad
 M_{\rm c}(t)=t^3+3\tau_{\rm c}G_{\rm c}t^2
                  -\tau_{\rm c}^2U_{\rm c}t+\tau_{\rm c}^3.
\]
The coefficients thus give exactly
\eqref{deg:compact-cubic-field}--\eqref{deg:compact-cubic-polynomials}.
Moreover
\(K_{\rm fin}=\tau_{\rm c}+O(\tau_{\rm c}^2)>0\), so the actual
finite multiplier tends to zero. The positive port images are buffered
by the fixed reserve above. The inward construction
\eqref{deg:compact-raw-tail-inward} and
\eqref{deg:compact-reflected-inward} gives separate actual positive
centre maps, with bounded nonlinear memory factors. Write their actions
as \(\mathcal A_{{\rm c},+}\) and \(\mathcal A_{{\rm c},-}\), based at
the original centre.

For \(u=a/\tau_{\rm c}\), direct integration across the bottleneck gives
\begin{equation}\label{rawstress:eq-cubic-positive-action}
 \mathcal A_{{\rm c},+}(u)
 =\frac{\pi_{\rm circ}B(1)}
             {(1+c_{\rm c})\kappa_{\rm c}}
                       +O(1+\tau_{\rm c}\log u).
\end{equation}
For completeness, put \(h(v)=vB(v)/(v+c_{\rm c})\). On a fixed
symmetric interval about \(v=1\), subtract \(h(1)+h'(1)(v-1)\).
The linear term integrates to zero against
\(((v-1)^2+\kappa_{\rm c}^2)^{-1}\), and the quadratic remainder
has bounded integral. Outside this interval the denominator has a
fixed margin; for large \(v\), \(h(v)=O(1+\tau_{\rm c}v)\), giving
the remaining \(O(1+\tau_{\rm c}\log u)\) bound. Thus
\(\kappa_{\rm c}\mathcal A_{{\rm c},+}\to\pi_{\rm circ}/2\).

At a fixed negative lip \(r_0<c_{\rm c}\), the negative action is
bounded. The positive centre height from
\eqref{rawstress:eq-cubic-positive-action} is smaller than its negative
centre height for small \(\tau_{\rm c}\). The actual image
\eqref{deg:compact-negative-centre-image} therefore gives a unique match,
before an observer is attached. Its radius tends to \(c_{\rm c}\):
on every fixed complement of that root, the negative action and
polynomial logarithm stay bounded, whereas the positive action diverges.
The negative root-collar construction in
Lemma~\ref{deg:lem-weighted-cubic-tail} gives a positive-time outward
passage on a fixed buffered invariant chord and continues it to the
original negative raw section. It applies for arbitrarily small \(X\).
This proves negative output and closing admission. Equality at the
same centre joins the two graphs by scalar uniqueness.

Both initial polynomial factors occur in the exact centre matching, so
\begin{equation}\label{rawstress:eq-cubic-memory}
 \mathcal A_{{\rm c},+}(u)-\mathcal A_{{\rm c},-}(r)
 =\tau_{\rm c}\log\frac{P_{\rm c}(u)}{P_{\rm c}(-r)}+O(\tau_{\rm c}).
\end{equation}
Put
\[
 L_-(r)=(r+1)^2+\kappa_{\rm c}^2,\qquad
 a_-:=\frac{c_{\rm c}B(-c_{\rm c})}{L_-(c_{\rm c})}\longrightarrow\frac14.
\]
The simple-root action and the source polynomial give, uniformly in a
fixed negative collar,
\[
 \mathcal A_{{\rm c},-}(r)=a_-\log(1/X)+O(1),\qquad
 \log\frac{P_{\rm c}(u)}{P_{\rm c}(-r)}
                               =3\log u+\log(1/X)+O(1).
\]
Inserting these in \eqref{rawstress:eq-cubic-memory} yields
\[
 (a_-+\tau_{\rm c})\log(1/X)
             =\mathcal A_{{\rm c},+}(u)
                              +O(1+\tau_{\rm c}\log u).
\]
This proves \eqref{rawstress:eq-cubic-root-distance}. Also
\(H=\ell_0M_{\rm c}(a)
e^{-\mathcal A_{{\rm c},+}(u)/\tau_{\rm c}}E_+\), with
\(E_+,E_+^{-1}\le C\) and \(M_{\rm c}(a)\to a^3\).
Equation~\eqref{rawstress:eq-cubic-positive-action} proves
\eqref{rawstress:eq-cubic-height}. The height is positive and below
the fixed cap. The raw speed at the centre is
\(\tau_{\rm c}^3+\ell(\gamma)H\), which tends to zero, so the
original speed restriction holds as well.

The root distance is too small to discard its power in the negative
output gain. We keep that power while forming the physical derivative.
With
\[
 \lambda_G=\frac{12\tau_{\rm c}G_{\rm c}}
                       {1+2\tau_{\rm c}G_{\rm c}},\qquad
 \beta_-=\frac{\tau_{\rm c}L_-(c_{\rm c})}
                    {c_{\rm c}B(-c_{\rm c})}=\frac{\tau_{\rm c}}{a_-},
\]
the complete collar gain retains
\begin{equation}\label{rawstress:eq-cubic-full-gain}
 \Gamma_{-,b_-}'=
 \tau_{\rm c}^{2-\lambda_G}\mathcal U_{\rm c}\,
 \tau_{\rm c}^{\beta_-}X^{\beta_-}
                      K_{\rm c}(X,X^{\beta_-}),\qquad b_-=-r^2/2.
\end{equation}
Here \(\mathcal U_{\rm c}\) and \(K_{\rm c}\) are bounded positive
units; the fixed output bridge and divided closing are included in
their bounded terminal factors. Equation
\eqref{rawstress:eq-cubic-root-distance} gives
\[
 \beta_-\log(1/X)\sim
             \frac{8\pi_{\rm circ}\tau_{\rm c}}{\kappa_{\rm c}}
                   \longrightarrow\infty,
 \qquad X^{\beta_-}\longrightarrow0.
\]
The input gain in \(b_+=-u^2/2\) satisfies
\(\Gamma_{+,b_+}'\asymp_{d_0}\tau_{\rm c}^2\) by
\eqref{deg:compact-raw-tail-input}. The physical quotient is exactly
\begin{equation}\label{rawstress:eq-cubic-physical-ratio}
 \Pi'=\frac{\Gamma_{-,b_-}'}{\Gamma_{+,b_+}'}
       \frac{B_+P_-}{B_-P_+}\frac{\Lambda_+}{\Lambda_-}.
\end{equation}
Since \(\lambda_G=O(\tau_{\rm c})\),
\(\tau_{\rm c}^{-\lambda_G}\) is bounded. The factors
\(\tau_{\rm c}^{\beta_-}\) and \(X^{\beta_-}\) are at most one.
The input polynomial satisfies \(P_+(u)\ge c(d_0)\tau_{\rm c}^{-3}\),
whereas \(P_-\le C\) and the two centre variation factors are units.
This proves
\(0<\Pi'\le C(d_0)\tau_{\rm c}^3\), with the original divided closing
and all exterior gains retained.
\end{proof}
